% -*-coding: utf-8 -*-

\zp{Fernande}{Georges Brassens}

\zs
<C>Une manie de vieux garçon
Moi <F>j'ai pris l'habi<E>tude

<G>D'agré<A7>menter ma<Dmi> soli<E>tude
Aux <D7>accents <G7>de cette chan<C>son <D7>
\ks

\zr 
Quand <G>je pense à Fer<Ami>nande
Je <D7>bande, je <G>bande

Quand j'pense à Féli<C>cie
Je bande <G>aussi

Quand j'pense à Lé<C>onore
Mon <D7>dieu je bande en<G>core

Mais <H7>quand j'pense à Lu<Emi>lu
<Ami>Là <D7>je ne bande <E>plus

La <H7>bandaison pa<Emi>pa
Ça n'se com<H7>man<D7>de <G>pas ~~~ <C>
\kr

\zs
C'est cette mâle ritournelle
Cette antienne virile

Qui retentit dans la guérite
De la vaillante sentinelle
\ks

\zr \kr

\zs
Afin de tromper son cafard
De voir la vie moins terne

Tout en veillant sur sa lanterne
Chante ainsi le gardien de phare
\ks

\zr \kr

\zs
Après la prière du soir
Comme il est un peu triste

Chante ainsi le séminariste
À genoux sur son reposoir
\ks

\zr \kr

\zs
À l'étoile où j'étais venu
Pour ranimer la flamme

J'entendis ému jusqu'aux larmes
La voix du soldat inconnu
\ks

\zr \kr

\zs
Et je vais mettre un point final
À ce chant salutaire

En suggérant au solitaire
D'en faire un hymne national
\ks

\zr \kr

\kp
